\documentclass[12pt,a4paper]{article}

\input{/home/tabaka/Documents/GitHub/Defs/defs.tex}

%\pagestyle{empty} 			%usuwa nr strony

\begin{document}
	
	\begin{center}
		\Large Całki nieoznaczone
	\end{center}\normalsize


	\ZbiorZadanieTwoXThree{Obliczyć poniższe całki nieoznaczone:}
	{$\displaystyle \int (x^4 - x^3 + x^2) \, dx$}
	{$\displaystyle \int (5t^8 - 2t^4 + t + 3) \, dt$}
	{$\displaystyle \int (7u^{3/2} + 2u^{1/2}) \, du$}
	{$\displaystyle \int (\frac{3}{x^{-2}} - \frac{4}{x^{-3}}) \, dx$}
	{$\displaystyle \int \left( \frac{4}{3t^2} + \frac{7}{2t} \right) dt$}
	{$\displaystyle \int \left( 5\sqrt{y} - \frac{3}{\sqrt{y}} \right) dy$}

	\vspace{0.4cm}

	\ZbiorZadanieTwoXTwo{Obliczyć poniższe całki nieoznaczone:}
	{$\displaystyle \int (2\sin \theta + 3\cos \theta) \, d\theta$}
	{$\displaystyle \int \cos^2 \frac{\theta}{2} \, d\theta$}
	{$\displaystyle \int \frac{d\theta}{\sin^2\theta \cos^2\theta}$}
	{$\displaystyle \int \tg^2\theta \, d\theta$}

	\vspace{0.4cm}
	
	\Wzor{$\displaystyle \int \frac{f'(x)}{f(x)} \, dx = \ln|f(x)|+C$}
	
	\ZbiorZadanieTwoXThree{Korzystając ze wzoru powyżej wyznaczyć całki nieocznaczone:}
	{$\displaystyle \int  \frac{2x}{1+x^2}\, dx$}
	{$\displaystyle \int  \frac{x}{3-5x^2}\, dx$}
	{$\displaystyle \int  \frac{4x-3}{2x^2-3x+5}\, dx$}
	{$\displaystyle \int  \ctg x\, dx$}
	{$\displaystyle \int  \frac{\cos2x}{\cos^2x\sin^2x}\, dx$}
	{$\displaystyle \int  \tg x\, dx$}

	\newpage

	\begin{center}
		\Large Całkowanie przez podstawianie
	\end{center}\normalsize

	\ZbiorZadanieTwoXThree{Obliczyć korzystając z całkowania przez podstawianie oraz obok podanego podstawienia:}
	{$\displaystyle \int e^{-2x+3}\, dx \qquad t=-2x+3$}
	{$\displaystyle \int \ctg(1+2x)\, dx \qquad t=1+2x$}
	{$\displaystyle \int xe^{-x^2}\, dx \qquad t=-x^2$}
	{$\displaystyle \int \frac{e^{\sqrt{x}}}{\sqrt{x}} \, dx \qquad t=\sqrt{x}$}
	{$\displaystyle \int \frac{x}{1+x^4}\, dx \qquad t=x^2$}
	{$\displaystyle \int \sqrt{x^3+1} x^2 \, dx \qquad t=x^3+1$}
	
	\vspace{0.4cm}
	
	\ZbiorZadanieTwoXThree{Obliczyć korzystając z całkowania przez podstawianie:}
	{$\displaystyle \int e^{x+3} + e^{x-3}\, dx$}
	{$\displaystyle \int (-3x+4)e^{-3x^2+8x}\, dx$}
	{$\displaystyle \int \frac{3x}{(x^2+1)^7}\, dx$}
	{$\displaystyle \int \frac{\sqrt{1+\sqrt{x}}}{\sqrt{x}}\, dx$}
	{$\displaystyle \int \frac{e^x}{e^{2x}+1}\, dx$}
	{$\displaystyle \int \frac{1}{x^2+9}\, dx \qquad \text{Hint:} t=\frac{1}{3}x$}
	
	\begin{center}
		\Large Całkowanie przez części
	\end{center}\normalsize
	
	\ZbiorZadanieTwoXThree{Obliczyć korzystając z całkowania przez części:}
	{$\displaystyle \int xe^x\, dx$}
	{$\displaystyle \int \cos^2x\, dx$}
	{$\displaystyle \int x^2e^{-3x}\, dx$}
	{$\displaystyle \int \ln x\, dx$}
	{$\displaystyle \int x\cos x\, dx$}
	{$\displaystyle \int e^x\sin x\, dx$}
	
	\vspace{0.4cm}
	
	\ZbiorZadanieTwoXThree{Obliczyć korzystając z całkowania przez części:}
	{$\displaystyle \int x\sin 10x\, dx$}
	{$\displaystyle \int x^2\cos x\, dx$}
	{$\displaystyle \int x\ln x\, dx$}
	{$\displaystyle \int x^5\ln x\, dx$}
	{$\displaystyle \int x2^x\, dx$}
	{$\displaystyle \int \arcsin x\, dx$}
	
	
	
\end{document}